\section{Discussion}
\label{sec:discussion}

In this section, we examine three properties of reflective code evolution: search diversity, seed stability, and category effects.

\paragraph{How does evolution explore the program space?}

Figure~\ref{fig:landscape}(c) provides another view of the search process.
Evolved programs form several regions with different storage designs, including LLM-centric, vector-only, SQL-only, SQL+vector, and vector+LLM programs.
Programs from different benchmarks often appear in the same regions.
This pattern suggests that evolution can try common storage designs before adapting them to a task.
Figure~\ref{fig:landscape}(a,b) further compares population-based search with a linear variant.
The linear variant starts from one empty seed and repeatedly mutates only the current best program\rev{; this corresponds to the $-$\,Diversity ablation in Table~\ref{tab:ablation}}.
Population search keeps multiple candidates active, so later mutations can draw from more than one candidate lineage.
The LoCoMo ablation matches this interpretation: removing diversity reduces test F1 from 0.459 to 0.318, and max sampling with multiple starters reaches 0.381 (Table~\ref{tab:ablation}).
These results suggest that population search helps preserve promising program lineages under a small iteration budget.

\paragraph{Is evolution stable across random seeds?}

% === Table: Stability Across Seeds ===
% Table: Stability Across Seeds
\begin{wraptable}{r}{0.64\textwidth}
    \vspace{-1.2em}
    \centering
    \scriptsize
    \caption{\textbf{Stability across evolution seeds.}
    Mean and std.\ are computed over five independent stability runs with random seeds; metrics match the primary columns in \autoref{tab:main-results}.
    CV is the coefficient of variation in percent, Best Baseline is the strongest fixed baseline, and Win is the number of seeds above that baseline.}
    \vspace{-0.5em}
    \label{tab:stability}
    \begin{tabularx}{\linewidth}{@{}l >{\centering\arraybackslash}X >{\centering\arraybackslash}p{0.6cm} >{\centering\arraybackslash}p{2.6cm} >{\centering\arraybackslash}p{0.6cm}@{}}
        \toprule
        \rowcolor{gray!10}
        \textbf{Benchmark} & \textbf{Mean{\scriptsize$\pm$Std}} & \textbf{CV} & \textbf{Best Baseline} & \textbf{Win} \\
        \midrule
        LoCoMo          & $\rev{0.421_{\pm0.042}}$ & \rev{10.0} & 0.373 {\scriptsize(Mem0)}         & 4/5 \\
        HB Data         & $0.391_{\pm0.023}$ &  5.9 & 0.327 {\scriptsize(GEPA{+}VS)}    & 5/5 \\
        PR Finance      & $0.561_{\pm0.032}$ &  5.7 & 0.449 {\scriptsize(GEPA)}          & 5/5 \\
        \bottomrule
    \end{tabularx}%
    \vspace{-1em}
\end{wraptable}


A practical concern for evolutionary methods is sensitivity to random seeds, especially under a 20-iteration budget in which discovery quality can depend on early mutations.
To evaluate this effect, Table~\ref{tab:stability} reports test scores from five independent runs on three benchmarks.
Across all benchmarks, the coefficient of variation remains \rev{at or below 10\%}, which indicates stable outcomes with moderate variance.
Across the 15 runs, 14 outperform the strongest baseline, supporting consistent gains across initializations.
Appendix~\ref{app:stability} provides full per-seed scores and validation trajectories.

\paragraph{Does evolution improve performance uniformly?}

% === Table: Per-Category Breakdown (combined) ===
% Table: Per-Category Breakdown (combined)
\begin{table*}[t]
    \centering
    \small
    \caption{\textbf{Per-category performance breakdown.}
    \textbf{(a)}~ALFWorld Unseen success rate by task type.
    \textbf{(b)}~LoCoMo token F1 by query type.
    $n$ is the number of test episodes or test queries per category.
    Min reports the minimum score across shown categories.
    Best per column in \textbf{bold}.}
    \vspace{-0.5em}
    \label{tab:per-category}
    \resizebox{\textwidth}{!}{%
    \begin{tabular}{@{}l ccccc c !{\vrule width 0.6pt} cccc c@{}}
        \toprule
        \rowcolor{gray!10}
        & \multicolumn{6}{c!{\vrule width 0.6pt}}{\textbf{(a) ALFWorld Unseen}} & \multicolumn{5}{c}{\textbf{(b) LoCoMo}} \\
        \cmidrule(r){2-7} \cmidrule(l){8-12}
        \rowcolor{gray!10}
        & \textbf{Simple} & \textbf{Clean} & \textbf{Cool} & \textbf{Heat} & \textbf{2-Obj} &
        & \textbf{Single} & \textbf{Temp.} & \textbf{Causal} & \textbf{Open} & \\
        \rowcolor{gray!10}
        & \footnotesize$(n{=}8)$ & \footnotesize$(n{=}11)$ & \footnotesize$(n{=}7)$ & \footnotesize$(n{=}6)$ & \footnotesize$(n{=}8)$ & \textbf{Min}
        & \footnotesize$(n{=}18)$ & \footnotesize$(n{=}18)$ & \footnotesize$(n{=}9)$ & \footnotesize$(n{=}55)$ & \textbf{Min} \\
        \midrule
        No Memory            & 0.88 & 0.82 & 0.57 & 0.83 & 0.62 & 0.57 & 0.02 & 0.00 & 0.11 & 0.04 & 0.00 \\
        Mem0                 & 0.88 & 0.64 & 0.57 & \textbf{1.00} & 0.75 & 0.57 & \textbf{0.33} & 0.26 & 0.37 & 0.43 & 0.26 \\
        Trajectory Retrieval & 0.75 & 0.73 & 0.71 & 0.67 & 0.75 & 0.67 & 0.18 & 0.33 & 0.34 & 0.28 & 0.18 \\
        GEPA + Vector Search & 0.88 & \textbf{1.00} & 0.57 & \textbf{1.00} & 0.62 & 0.57 & 0.15 & \textbf{0.38} & \textbf{0.39} & 0.31 & 0.15 \\
        \rowcolor{green!10}
        \sysname{}      & 0.88 & 0.91 & \textbf{1.00} & 0.83 & 0.75 & \textbf{0.75} & \textbf{0.33} & 0.19 & 0.35 & \textbf{0.51} & 0.19 \\
        \bottomrule
    \end{tabular}}%
    \vspace{-1em}
\end{table*}


Table~\ref{tab:per-category} disaggregates the results in Table~\ref{tab:main-results} by category (full per-category tables are in Appendix~\ref{app:per-category}).
On ALFWorld Unseen, \sysname{} raises the worst category score to 0.75, above the best baseline floor among the methods shown (0.67).
This pattern suggests that the evolved action cache addresses failure modes that fixed designs leave exposed across task types.
On LoCoMo, the gains are less uniform: \sysname{} is strongest on open-domain questions and tied on single-hop recall, while Mem0 and GEPA+Vector Search remain stronger on temporal and causal questions.
This split suggests that an aggregate objective can improve the overall score while leaving lower-frequency or harder categories underserved.
